\documentclass[11pt]{article}
\usepackage[margin=1in]{geometry}
\usepackage{enumitem}
% =============================================================================
% Preambles/header.tex — shared LaTeX/Beamer preamble for this project
%
% Use from a Beamer lecture by prepending:
%     \input{header}
% and compiling with TEXINPUTS=../Preambles:$TEXINPUTS (the /compile-latex
% skill does this automatically).
%
% PALETTE CONTRACT. Color names defined here MUST match the names in
% Quarto/theme-template.scss so Beamer and Quarto renderings use the same
% palette. The scripts/check-palette-sync.sh script enforces this.
%
% To customize: edit the HEX values below AND the matching $variable in
% Quarto/theme-template.scss, then run scripts/check-palette-sync.sh.
% =============================================================================

% -----------------------------------------------------------------------------
% Engine detection + encoding
%
% Our /compile-latex skill uses XeLaTeX, where inputenc is unnecessary and
% can produce encoding warnings. Only load inputenc under pdfTeX.
% -----------------------------------------------------------------------------
\usepackage{iftex}
\ifPDFTeX
  \usepackage[utf8]{inputenc}
\fi

\usepackage{amsmath, amssymb, amsthm}
\usepackage{booktabs}
\usepackage{graphicx}

% -----------------------------------------------------------------------------
% PALETTE — mirrors Quarto/theme-template.scss variable names.
% Load xcolor and define colors BEFORE anything that references them
% (notably hyperref, hypersetup, and the Beamer color assignments).
% -----------------------------------------------------------------------------
\usepackage{xcolor}

\definecolor{primary-blue}{HTML}{012169}      % headings, accents
\definecolor{primary-gold}{HTML}{B9975B}      % emphasis, borders
\definecolor{highlight-yellow}{HTML}{F2A900}  % markers, alerts
\definecolor{light-bg}{HTML}{E8EDF5}          % title / section backgrounds
\definecolor{jet}{HTML}{1A1A1A}               % body text

% Semantic colors (used in diagrams and annotations)
\definecolor{positive}{HTML}{15803D}          % good / observed / identified
\definecolor{negative}{HTML}{B91C1C}          % bad / problematic / bias
\definecolor{neutral}{HTML}{525252}           % reference / context

% Highlight accents (match .hi-* classes in the SCSS)
\definecolor{hi-slate}{HTML}{314F4F}
\definecolor{hi-green}{HTML}{2E8B57}
\definecolor{hi-red}{HTML}{C41E3A}

% Semantic aliases so skill templates and snippets can write
% \textcolor{accent}{...} without hard-coding "primary-blue".
\colorlet{accent}{primary-blue}
\colorlet{accent2}{primary-gold}

% -----------------------------------------------------------------------------
% Hyperref — loaded AFTER xcolor + palette so link colors resolve.
% -----------------------------------------------------------------------------
\usepackage{hyperref}
\hypersetup{colorlinks=true, linkcolor=primary-blue, urlcolor=primary-blue,
            citecolor=primary-blue, pdfborder={0 0 0}}

% -----------------------------------------------------------------------------
% Course code — set once per deck (after \input{header}, before \begin{document})
% via \coursecode{CS401}. Shown in the footer on every slide so a deck's course
% is identifiable without opening the title page. Empty by default.
% -----------------------------------------------------------------------------
\makeatletter
\newcommand{\coursecode}[1]{\gdef\@coursecodetext{#1}}
\gdef\@coursecodetext{}
\makeatother

% -----------------------------------------------------------------------------
% Beamer theme assignments (only apply under Beamer).
%
% We detect Beamer via \@ifclassloaded{beamer}, which is the documented,
% non-internal way. This must be wrapped in \makeatletter/\makeatother
% because \@ifclassloaded contains an @-catcode character.
% -----------------------------------------------------------------------------
\makeatletter
\@ifclassloaded{beamer}{%
  \usefonttheme{professionalfonts}
  \setbeamercolor{structure}{fg=primary-blue}
  \setbeamercolor{frametitle}{fg=primary-blue}
  \setbeamercolor{title}{fg=primary-blue}
  \setbeamercolor{subtitle}{fg=primary-gold}
  \setbeamercolor{author}{fg=primary-blue}
  \setbeamercolor{institute}{fg=neutral}
  \setbeamercolor{date}{fg=primary-gold}
  \setbeamercolor{itemize item}{fg=primary-blue}
  \setbeamercolor{itemize subitem}{fg=primary-gold}
  \setbeamercolor{alerted text}{fg=primary-gold}
  \setbeamercolor{block title}{fg=white, bg=primary-blue}
  \setbeamercolor{block body}{bg=light-bg}
  % Semantic block variants: block=definition/concept (blue), exampleblock=
  % worked example (green/positive), alertblock=key takeaway or caution (gold).
  % Keep to <=2 colored boxes per slide across ALL three variants (INV-7).
  \setbeamercolor{block title example}{fg=white, bg=positive}
  \setbeamercolor{block body example}{bg=light-bg}
  \setbeamercolor{block title alerted}{fg=white, bg=primary-gold}
  \setbeamercolor{block body alerted}{bg=light-bg}

  % Minimal footer: course code (left) + page X / N (right), no navigation clutter
  \setbeamertemplate{navigation symbols}{}
  \setbeamertemplate{footline}{%
    \color{neutral}\footnotesize\hspace{0.7em}\@coursecodetext\hfill
    \insertframenumber/\inserttotalframenumber\hspace{0.7em}\vspace{0.3em}}
}{}
\makeatother

% -----------------------------------------------------------------------------
% TikZ setup — libraries the snippet gallery uses
% -----------------------------------------------------------------------------
\usepackage{tikz}
\usetikzlibrary{arrows.meta, positioning, calc, decorations.pathreplacing,
                fit, shapes.geometric, backgrounds}

% Shared TikZ styles — snippets and hand-written diagrams can pick these up
% via `\begin{tikzpicture}[every node/.style={...}]` or by naming specific
% styles (e.g., `\node[dag-node] (X) at (0,0) {X};`).
\tikzset{
  % Node styles for causal DAGs and similar
  dag-node/.style = {
    draw=primary-blue, fill=light-bg, rounded corners=2pt,
    minimum width=1.2cm, minimum height=0.9cm, align=center, font=\small
  },
  decision-node/.style = {
    draw=primary-gold, fill=white, diamond, aspect=1.8,
    minimum width=1.8cm, minimum height=1.2cm, align=center, font=\small,
    inner sep=1pt
  },
  % Edge styles — observed vs counterfactual
  observed-edge/.style = {-{Latex[length=2.5mm]}, thick, draw=jet},
  counterfactual-edge/.style = {-{Latex[length=2.5mm]}, thick, draw=neutral, dashed},
  confound-edge/.style = {-{Latex[length=2.5mm]}, thick, draw=negative, dashed},
  % Data-point styles for plots
  observed-dot/.style = {fill=primary-blue, draw=primary-blue, circle, inner sep=1.2pt},
  counterfactual-dot/.style = {fill=white, draw=neutral, circle, inner sep=1.2pt}
}

% -----------------------------------------------------------------------------
% Convenience macros
% -----------------------------------------------------------------------------
\newcommand{\muted}[1]{\textcolor{neutral}{#1}}
\newcommand{\key}[1]{\textcolor{primary-gold}{\textbf{#1}}}
\newcommand{\good}[1]{\textcolor{positive}{#1}}
\newcommand{\bad}[1]{\textcolor{negative}{#1}}

% Standout frame macro: full-bleed dark background for section transitions.
% Beamer-only (uses \begin{frame}/\setbeamercolor/\usebeamerfont) -- do not
% call from an article-class file (e.g. Notes/<CODE>/*-notes.tex). \muted,
% \key, \good, \bad, and \coursecode above are plain xcolor/gdef and are
% safe to use from either class; verified when Notes/article-class support
% was added.
% Expand: \transitionslide{Your transition title}
\newcommand{\transitionslide}[1]{%
  \begingroup
  \setbeamercolor{background canvas}{bg=primary-blue}
  \begin{frame}[plain]
    \centering
    \vfill
    {\usebeamerfont{title}\color{white}\Large #1\par}
    \vfill
  \end{frame}
  \endgroup
}

% Section divider frame (Beamer-only), an alternative to \transitionslide with a
% darker two-line style: near-black (jet) background, a small white label line
% above a larger gold title, and a thin gold rule along the bottom edge.
% Expand: \sectiondivider{Section 2}{Pointers}
\newcommand{\sectiondivider}[2]{%
  \begingroup
  \setbeamercolor{background canvas}{bg=jet}
  \begin{frame}[plain]
    \vspace*{3.0cm}
    {\color{white}\ttfamily\large #1\par}
    \vspace{0.5cm}
    {\color{primary-gold}\LARGE #2\par}
    \begin{tikzpicture}[remember picture, overlay]
      \draw[primary-gold, line width=2.5pt]
        ($(current page.south west)+(0,0.1)$) -- ($(current page.south east)+(0,0.1)$);
    \end{tikzpicture}
  \end{frame}
  \endgroup
}

% -----------------------------------------------------------------------------
% Channel watermark — "Concepts That Click" (YouTube).
%
% Article-class files (CompetitiveExam Q/A sets, Notes, Assignments) get a
% light diagonal draft watermark on every page plus a centered footer naming
% the channel. Beamer decks get a subtle TikZ background watermark instead
% (the footline already carries the course code). Centralized here so every
% MCQ PDF is branded without per-file edits.
% -----------------------------------------------------------------------------
\makeatletter
\@ifclassloaded{article}{%
  \usepackage{draftwatermark}
  \SetWatermarkText{Concepts That Click}
  \SetWatermarkScale{1}
  \SetWatermarkFontSize{2cm}
  \SetWatermarkLightness{0.86}
  \SetWatermarkAngle{45}
  \usepackage{fancyhdr}
  \pagestyle{fancy}
  \fancyhf{}
  \fancyfoot[C]{\footnotesize\textcolor{neutral}{Concepts That Click~|~YouTube: \href{https://www.youtube.com/@concepts.thatclick}{@concepts.thatclick}}}
  \renewcommand{\headrulewidth}{0pt}
  \renewcommand{\footrulewidth}{0pt}
  \fancypagestyle{plain}{%
    \fancyhf{}%
    \fancyfoot[C]{\footnotesize\textcolor{neutral}{Concepts That Click~|~YouTube: \href{https://www.youtube.com/@concepts.thatclick}{@concepts.thatclick}}}%
    \renewcommand{\headrulewidth}{0pt}%
    \renewcommand{\footrulewidth}{0pt}%
  }
}{}
\@ifclassloaded{beamer}{%
  \setbeamertemplate{background}{%
    \begin{tikzpicture}[remember picture, overlay]
      \node[rotate=45, opacity=0.055, align=center]
        at (current page.center)
        {\Huge\textcolor{neutral}{Concepts That Click}};
    \end{tikzpicture}%
  }
}{}
\makeatother 
\coursecode{JKSSB}
\renewcommand{\thesection}{62.\arabic{section}}
\newtheorem{example}{Example}[section]
\newenvironment{definitionbox}[1]{%
  \par\noindent\textbf{Definition (#1).}\ \itshape
}{\par\vspace{0.3em}}
\title{Current Internet, Networking and Cloud IaaS --- Detailed Lecture Notes}
\author{JKSSB Computer Awareness}
\date{\today}

\begin{document}
\maketitle

\section{TCP: ordered and reliable, but with no timing promise}
A running Internet conversation must survive lost, duplicated, and reordered
packets. \textbf{TCP} (Transmission Control Protocol) is the transport
protocol that hides this mess from the application. It is
\emph{connection-oriented}: the two ends set up a connection before any data
moves. Once the connection exists, TCP numbers every byte, acknowledges what
arrives, and retransmits whatever is lost, so the receiving application sees
the bytes in exactly the order they were sent. That is the double guarantee
to memorise: \textbf{ordered} delivery and \textbf{reliable} delivery.

What TCP does not promise is timing. It never claims the data will arrive
within a fixed time, because retransmissions take as long as they take. This
``no timing guarantee'' point is tested directly (PYQ-18): an option that
claims TCP promises delivery by a deadline is wrong.

\begin{definitionbox}{TCP}
A connection-oriented transport protocol that delivers bytes in order and
retransmits lost segments, with no guarantee about delivery time.
\end{definitionbox}

\begin{center}
\begin{tabular}{p{0.24\textwidth}p{0.66\textwidth}}
\toprule
\textbf{Term} & \textbf{Meaning in this lesson} \\
\midrule
TCP & Transmission Control Protocol; ordered and reliable transport. \\
UDP & User Datagram Protocol; connectionless, fast, unreliable transport. \\
ICMP & Internet Control Message Protocol; control and diagnostic messages only. \\
BGP & Border Gateway Protocol; inter-domain routing between autonomous systems. \\
IPv6 & 128-bit next-generation Internet addressing; no broadcast. \\
QUIC & Quick UDP Internet Connections; the UDP-based transport under HTTP/3. \\
DNS & Domain Name System; translates domain names into IP addresses. \\
PIR & Public Interest Registry; operates the .org registry. \\
ICANN & Internet Corporation for Assigned Names and Numbers; global name/number coordination. \\
ARPANET & Advanced Research Projects Agency Network (1969); ancestor of the Internet. \\
IaaS & Infrastructure as a Service; provider manages hardware, customer manages OS and above. \\
\bottomrule
\end{tabular}
\end{center}

The rest of this lesson uses these terms in one consistent sense. \textbf{TCP}
means ordered and reliable transport. \textbf{UDP} means the fast,
connectionless alternative with neither guarantee. \textbf{ICMP} means the
control-message protocol that never carries application data.
\textbf{BGP} means routing between autonomous systems, never inside one LAN.
\textbf{IPv6} means 128-bit addresses with no broadcast. \textbf{QUIC} means
the UDP-based transport that HTTP/3 uses instead of TCP. \textbf{DNS} means
the name-to-address directory. \textbf{PIR} runs the .org registry under
\textbf{ICANN} coordination. \textbf{ARPANET} is the 1969 packet-switched
ancestor of the Internet. \textbf{IaaS} splits responsibility so that the
provider manages networking hardware while the customer manages the operating
system, applications, and data.

\section{UDP and the TCP/UDP contrast}
\textbf{UDP} (User Datagram Protocol) is TCP's near-neighbour and the most
common distractor for it. UDP is \emph{connectionless}: each datagram is sent
independently with no setup. It provides no ordering guarantee, no
reliability guarantee, and no retransmission. Its header is small, so it is
faster and cheaper per packet than TCP.

\begin{definitionbox}{UDP}
A connectionless transport protocol that sends independent datagrams with no
ordering, reliability, or retransmission guarantee.
\end{definitionbox}

\begin{center}
\begin{tabular}{p{0.20\textwidth}p{0.35\textwidth}p{0.35\textwidth}}
\toprule
\textbf{Point} & \textbf{TCP} & \textbf{UDP} \\
\midrule
Connection & Connection-oriented & Connectionless \\
Ordering & Guaranteed & Not guaranteed \\
Reliability & Retransmits lost data & No retransmission \\
Speed & Slower (acknowledgements, state) & Faster (light header) \\
Typical use & Web pages, mail, file transfer & Live voice, video, games, DNS lookups \\
\bottomrule
\end{tabular}
\end{center}

The exam tests this pair by swapping one property: an option that says UDP
guarantees order, or that TCP is connectionless, is always wrong. The anchor
is the direction of the trade: TCP buys correctness at the cost of speed,
while UDP buys speed at the cost of correctness.

\begin{example}
An item states, ``UDP guarantees ordered delivery of datagrams.'' This is
wrong. Ordering is a TCP guarantee; UDP provides no ordering at all. The
wrong option has moved a TCP property onto UDP.
\end{example}

\section{ICMP: the network's own messenger}
While TCP and UDP carry user data, \textbf{ICMP} (Internet Control Message
Protocol) carries messages \emph{about} the network. Routers and hosts use it
to report errors --- destination unreachable, time exceeded, parameter
problems --- and diagnostic tools such as \textbf{ping} (echo request and
reply) and \textbf{traceroute} are built on ICMP messages.

\begin{definitionbox}{ICMP}
The control and diagnostic protocol of the Internet layer; reports errors and
supports ping and traceroute, and never carries application data.
\end{definitionbox}

Three facts together close every ICMP trap (PYQ-18, PYQ-19, TRAP-19). First,
ICMP transfers no application data: it never moves a file, a mail message,
or a web page. Second, it uses no ports: ports belong to TCP and UDP, and an
option that assigns ICMP a port number is wrong. Third, it is neither a
routing protocol like BGP nor a mail protocol like SMTP; it only reports on
conditions the other protocols encounter.

\section{BGP: routing between autonomous systems}
The Internet is not one network but a federation of large independently
managed networks called \textbf{autonomous systems} (an ISP, a university
backbone, a government network). \textbf{BGP} (Border Gateway Protocol) is
the protocol that routes \emph{between} them. It exchanges reachability
information across autonomous-system borders and selects the paths that
traffic takes over the Internet backbone.

\begin{definitionbox}{BGP}
The inter-domain routing protocol that exchanges reachability between
autonomous systems on the Internet backbone.
\end{definitionbox}

The tested contrast is BGP against the \emph{interior} protocols
\textbf{OSPF} (Open Shortest Path First) and \textbf{RIP} (Routing
Information Protocol), which route \emph{inside} a single domain (PYQ-19,
TRAP-20). Any option that places BGP inside one small local network has
swapped the two scopes. The one-word shortcut is ``inter-domain'': whenever
the stem says inter-domain or between autonomous systems, the answer is BGP.

\section{IPv6: address space, header, and no broadcast}
\textbf{IPv4} uses 32-bit addresses (about 4.3 billion), which ran out under
the growth of the Internet. \textbf{IPv6} replaces them with \textbf{128-bit}
addresses, written as eight groups of hexadecimal digits, giving about
$3.4 \times 10^{38}$ addresses. This vast space is the standard reason given
for IPv6 removing the need for address-sharing workarounds such as
\textbf{NAT} (Network Address Translation), which lets many private addresses
share one public IPv4 address (PYQ-19).

\begin{definitionbox}{IPv6 address}
A 128-bit Internet-layer address, written as eight hexadecimal groups, with a
fixed 40-byte header.
\end{definitionbox}

Two further IPv6 facts are tested. The fixed IPv6 header is \textbf{40 bytes}
long (PYQ-18). And IPv6 has \textbf{no broadcast}: IPv4's three delivery
styles (unicast to one receiver, multicast to a group, broadcast to everyone)
become unicast, multicast, and anycast (nearest of several receivers) in
IPv6, with a multicast to the all-nodes group replacing the broadcast. An
option mentioning an ``IPv6 broadcast address'' is therefore self-defeating.
The 256-bit figure that appears as a distractor is always wrong (TRAP-18).

\begin{example}
An item offers ``IPv6 uses 256-bit addresses'' as one option. This is wrong.
IPv6 uses 128-bit addresses; 256 is double the correct size and matches no
Internet-layer standard. The option confuses IPv6 with an unrelated larger
number.
\end{example}

\section{HTTP/3 and QUIC}
Web traffic travelled over TCP for decades: HTTP/1.1 and HTTP/2 both run
over TCP. \textbf{HTTP/3} breaks that pattern. It runs over \textbf{QUIC}
(Quick UDP Internet Connections), which itself runs over \textbf{UDP}, not
over TCP. QUIC rebuilds the needed reliability, congestion control, and
security on top of UDP, gaining faster connection setup and independent
streams without head-of-line blocking across streams.

\begin{definitionbox}{QUIC}
A UDP-based transport protocol carrying HTTP/3, providing its own reliability
and security instead of using TCP.
\end{definitionbox}

The trap writes itself: any option claiming HTTP/3 relies on TCP has the
protocol stack backwards (PYQ-19, TRAP-21). The correct chain is HTTP/3 over
QUIC over UDP.

\section{DNS and the .org / PIR / ICANN chain}
\textbf{DNS} (Domain Name System) is the Internet's directory: it translates
human-readable domain names such as \texttt{example.org} into numeric IP
addresses that routers can use. Resolution runs hierarchically from root
servers down through top-level-domain servers to the domain's own
authoritative servers. DNS is a lookup service, not a routing protocol and
not a mail-transfer protocol, and confusing it with either is a standard
wrong answer.

The governance half of the topic names three actors. \textbf{.org} is a
generic top-level domain originally meant for non-profit organisations.
Its registry is operated by \textbf{PIR} (Public Interest Registry).
\textbf{ICANN} (Internet Corporation for Assigned Names and Numbers)
coordinates domain names and IP addresses at the global level. The exam
distinction is role-based: PIR \emph{runs} the .org registry, while ICANN
\emph{coordinates} the naming system as a whole.

\section{ARPANET and packet switching}
\textbf{ARPANET} (Advanced Research Projects Agency Network), live in
\textbf{1969}, is the ancestor of the Internet. Its contribution was proving
that \textbf{packet switching} works across long distances: splitting data
into addressed packets that travel independently and are reassembled at the
destination, with the network routing around failed links. Later ARPANET work
produced the TCP/IP protocol family.

\begin{definitionbox}{Packet switching}
Splitting data into addressed packets that may take different routes and are
reassembled at the destination.
\end{definitionbox}

The tested contrast is packet switching against \textbf{circuit switching},
the older telephone model that holds one dedicated path for the whole call.
The Internet uses packet switching; the traditional telephone call used
circuit switching. Options that assign a dedicated reserved circuit to the
Internet have the two models swapped.

\section{The router: joining networks by reading IP addresses}
A \textbf{router} connects different networks and forwards packets between
them. It does this by reading each packet's destination \textbf{IP address}
and choosing the next hop from its \textbf{routing table}. In
assertion--reason form: the assertion (a router joins networks) and the
reason (it reads the destination IP and consults its table) are both true,
and the reason correctly explains the assertion, because address-reading is
\emph{how} the router forwards between networks.

The near-neighbours are the \textbf{switch}, which forwards \emph{inside} one
network using MAC addresses, and the \textbf{hub}, which merely repeats
everything it receives. An option that says a router forwards within one
network using MAC addresses describes a switch, not a router.

\section{IaaS: the provider manages the hardware}
Cloud service models split responsibility between provider and customer.
In \textbf{IaaS} (Infrastructure as a Service), the \textbf{provider} manages
the physical infrastructure: servers, storage hardware, \textbf{networking
hardware}, and virtualisation. The \textbf{customer} manages everything above
that layer: the operating system, runtime, applications, and data. The exam
anchor is the provider's side: in IaaS the provider manages the networking
hardware (PYQ-20).

\begin{definitionbox}{IaaS responsibility split}
The provider manages infrastructure (networking and storage hardware,
virtualisation); the customer manages the operating system, runtime,
applications, and data.
\end{definitionbox}

Moving down the ladder, the provider manages more. In \textbf{PaaS}
(Platform as a Service) the provider additionally manages the operating
system and runtime, leaving the customer its applications and data. In
\textbf{SaaS} (Software as a Service) the provider manages everything up to
the finished application, and the customer only manages its own data and
settings. The trap claims the IaaS provider manages the customer's OS or
data; it never does (TRAP-22).

\begin{example}
An item states, ``In IaaS the provider manages the customer's operating
system and data.'' This is wrong. Those belong to the customer; the provider
stops at the infrastructure layer (networking and storage hardware plus
virtualisation). The option describes SaaS, not IaaS.
\end{example}

\section{Exam retrieval and common confusions}
Seven distinctions prevent most errors on this topic.
\begin{enumerate}
  \item TCP guarantees \textbf{order and reliability} but no timing; UDP
    guarantees \textbf{neither} (PYQ-18).
  \item ICMP carries \textbf{control messages only}: no application data, no
    ports (TRAP-19).
  \item BGP is \textbf{inter-domain} (between autonomous systems); OSPF and
    RIP are intra-domain (TRAP-20).
  \item IPv6 is \textbf{128-bit} with a \textbf{40-byte} header; 256-bit is
    always wrong (PYQ-18, TRAP-18).
  \item IPv6 has \textbf{no broadcast}; the replacement is
    \textbf{multicast} to the all-nodes group (PYQ-19).
  \item HTTP/3 uses \textbf{QUIC over UDP}, never TCP (TRAP-21).
  \item In IaaS the \textbf{provider} manages networking hardware; the
    \textbf{customer} manages the OS, applications, and data (TRAP-22).
\end{enumerate}

\begin{example}
An item asks which protocol routes between autonomous systems on the
Internet backbone. Trace the scope word: ``between autonomous systems'' is
the definition of inter-domain routing, and the inter-domain protocol is BGP.
OSPF and RIP are interior protocols, and ICMP is not a routing protocol at
all, so all three are eliminated by scope.
\end{example}

\subsection*{Exercises}
\begin{enumerate}[label=\arabic*.]
  \item State the two guarantees TCP provides and the one it does not.
  \item Distinguish TCP from UDP on connection, ordering, reliability, and
    speed.
  \item What does ICMP carry, and what two things does it never do?
  \item Distinguish BGP from OSPF/RIP by scope.
  \item State the IPv6 address size and header size, and explain what replaces
    the broadcast.
  \item Explain the HTTP/3 protocol stack in one chain of three names.
  \item Distinguish the roles of PIR and ICANN for the .org domain.
  \item State the IaaS responsibility split in two sentences.
\end{enumerate}

\subsection*{Solutions}
\begingroup\small
\begin{enumerate}[label=\arabic*.]
  \item TCP guarantees ordered delivery and reliable delivery (lost segments
    are retransmitted). It provides no timing guarantee.
  \item TCP is connection-oriented, ordered, and reliable but slower; UDP is
    connectionless, unordered, and unreliable (no retransmission) but faster.
  \item ICMP carries control and diagnostic messages (errors, ping,
    traceroute). It never transfers application data and never uses ports.
  \item BGP routes between autonomous systems (inter-domain, Internet
    backbone); OSPF and RIP route inside a single domain.
  \item IPv6 uses 128-bit addresses with a fixed 40-byte header. It has no
    broadcast; a multicast to the all-nodes group replaces it.
  \item HTTP/3 runs over QUIC, which runs over UDP: HTTP/3 over QUIC over
    UDP.
  \item PIR operates the .org registry; ICANN coordinates domain names and
    IP addresses globally.
  \item The provider manages the infrastructure (networking and storage
    hardware, virtualisation); the customer manages the operating system,
    runtime, applications, and data.
\end{enumerate}
\endgroup
\end{document}
